\label{ch_FunctionSecretSharing}
\subsection{(k,n) Threshold Schemes / Secret Sharing Schemes}
\label{sub:ThresholdSchemes}
Liu \cite{liu1968combMath} describes the general problem as follows: 
\begin{quote}
Eleven scientists are working on a secret project. 
They wish to lock up the documents in a cabinet so that the cabinet can be opened if and only if six or more of the scientists are present.
What is the smallest number of locks needed?
What is the smallest number of keys to the locks each scientist must carry?
\end{quote}
The minimal solution uses 462 locks and 252 keys per scientist \cite{shamir1979share}.
These numbers become exponentially worse when the number of scientists increases.

Secret Sharing is the generalization of this problem \cite{brickell1989schemes}.
A secret $D$ is divided into $n$ shares $D_1,...,D_n$ such that:
\begin{itemize}
    \item At least $k\leq n$ shares are necessarry to make $D$ computable
    \item If less than $k$ shares are present, $D$ cannot be efficiently computed
\end{itemize} 
Shamir \cite{shamir1979share} calls this a \textit{(k, n) Threshold Scheme}.
Boyle et al. \cite{boyle2015function} describe the same phenomenon as a \textit{Secret Sharing Scheme}. 
This research uses the terminology of Boyle et al.

One possible implementation by Shamir \cite{shamir1979share} is based on polynomial interpolation:
$D$ is represented by a random polynomial $q(x)=a_0+a_1x+...+a_{k-1}x^{k-1}$ with $a_0=D$ and degree $k-1$.
Each part is then represented by $D_1=q(1),...,D_i=q(i),...,D_n=q(n)$.
Therefore, one can only calculate $D$ if she has knowledge of at least $k-1$ of these values.  

\subsection{Additive Secret Sharing Scheme}
\label{sub:AdditiveSecretSharingScheme}
Boyle et al. \cite{boyle2015function} consider additive secret sharing the simplest type of secret sharing.

The secret $s$ is an element of an Abelian Group $G$.
The secret $s$ is reconstructed by adding all $p$ shares $s=\sum_{i=0}^p s_i$ (in $G$). 
Every subset of $p-1$ shares reveals nothing about the secret.
This implies every share $s_b$ on its own hides the secret $s$.

\subsection{Function Secret Sharing Scheme}
\label{sub:FunctionSecretSharingScheme}
A \textit{Function Secret Sharing Scheme (FSS)} is the natural extension of an additive secret sharing scheme \cite{boyle2016function}.
Thus, high-level correctness requirements stay the same: A secret is shared among $m$ parties such that no subset of less than $k$ participants can reconstruct it.
In this case, the secret $f\in F$ is a function $f:\{0,1\}^n\rightarrow G$ with $G$ being an abelian group and $F$ being the family of functions of type $f$ (\cite{boyle2015function}).

Boyle et al. \cite{boyle2015function} define such scheme by a pair of algorithms ($Gen$, $Eval$).
$Gen$ requires the definition of a function $f\in F$.
It will then output $m$ keys $(k_1,...,k_m)$ with each key $k_i$ defining the function $f_i(x)=Eval(i,k_i,x)$.
An input $x\in \{0,1\}^n$ is then shared among all participants so everyone evaluates on the same input.
With addition in $G$ the secret is evaluated via $f(x)=f_1(x)+...+f_m(x)$.
Again, every strict subset of size smaller than $k$ computationally hides $f$ (\cite{boyle2016function}).

An efficient FSS implementation is \textit{Distributed Point Functions} (\ref{sub:DistributedPointFunctions}).